
% Start preamble
\documentclass[12pt,a4paper]{article}
\usepackage{geometry}
 \geometry{
 a4paper,
 total={170mm,257mm},
 left=20mm,
 top=20mm,
 }
\usepackage[utf8]{inputenc}
\usepackage[T1]{fontenc}
\usepackage[pdftex]{graphicx}
\graphicspath{{./}}
\usepackage{enumitem}
\usepackage{pdfpages}
\usepackage{hyperref}
\usepackage{tikz}
\usepackage{attachfile}
\usepackage{epstopdf}
\usepackage{array}
\usepackage{multirow}
\usepackage{multicol}
\usepackage{float}
%\usepackage[table]{xcolor,colorbl}
\setlength{\textwidth}{17cm}
\setlength{\oddsidemargin}{-0.5cm}
\setlength{\evensidemargin}{-0.5cm}
%\setlenght{\headsep}{0cm}
\setlength\parindent{0pt}
%\setlength{\extrarowheight}{3pt}
\usepackage{listings}
%\usepackage{xcolor}

\input{arduinoLanguage.tex}
%%%%%% Counting oppgaves %%%%%%
 \newcount\questnum \questnum=0
 \def\oppgave{
            \advance\questnum by 1
	    \ifthenelse{\questnum=1}
	    {
                \vskip 1cm
		\textbf{Oppgave}\hskip 5pt\the\questnum \hfill \hfill(12p)
		\vskip 3pt
		\hrule
	\vskip 0.5cm}
	{
                \vskip 1cm
		\textbf{Oppgave}\hskip 5pt \the\questnum \hfill \hfill(6p)
		\vskip 3pt \hrule \vskip 0.5cm }

		}

% End preamble

\begin{document}
\title{Prøve Datakommuniasjon 01}
\author{Faglærer: Fred-Olav Mosdal 90507684\\
Oppgave 1-8 gjøres på ark. Hjelpemidler: kalkulator. \\
Oppgave 9 gjøres på PC. Hjelpemidler: Alle ikke kommuniserende\\
Elektronisk del av prøve sendes på email til: \\
fred-olav.mosdal@skole.rogfk.no\\
Emne: Høstprøve}
\maketitle
\oppgave{}
Basert på de korte beskrivelsene fra teknisk sjef blir det en diskusjon om hvordan de ulike enhetene skal kommunisere med hverandre. Hver av dere blir bedt om å foreslå en løsning som et grunnlag for diskusjon.
\vskip 5pt 
• Lag en skisse over anleggets kommunikasjonstopologi og begrunn valgene
dine
\vskip 2.5pt 
\newpage
\includepdf[pages=1,angle=0]{../output/nogpl/aKom01.pdf}
\oppgave{}
Hva menes med MAC-adresse? 
\vskip 2.5pt 
\begin{tikzpicture}
	\draw[step=0.5cm,gray!20,very thin]  grid (17,3) ;
\end{tikzpicture}
\vskip 2.5pt 
\oppgave{}
Hva står WAN, LAN og LAN  
\vskip 2.5pt 
\begin{tikzpicture}
	\draw[step=0.5cm,gray!20,very thin]  grid (17,3) ;
\end{tikzpicture}
\vskip 2.5pt 
\oppgave{}
Hva er en modulærtang? 
\vskip 2.5pt 
\begin{tikzpicture}
	\draw[step=0.5cm,gray!20,very thin]  grid (17,3) ;
\end{tikzpicture}
\vskip 2.5pt 
\oppgave{}
Hva betyr følgende merking på en TP-kabel SF/UTP 
\vskip 2.5pt 
\begin{tikzpicture}
	\draw[step=0.5cm,gray!20,very thin]  grid (17,3) ;
\end{tikzpicture}
\vskip 2.5pt 
\newpage
\oppgave{}
Tegn og forklar virkemåten til en switch 
\vskip 2.5pt 
\begin{tikzpicture}
	\draw[step=0.5cm,gray!20,very thin]  grid (17,15) ;
\end{tikzpicture}
\vskip 2.5pt 
\oppgave{}
Hvilke virkemidler kan en bruke for å øke datasikkerheten.  
\vskip 2.5pt 
\begin{tikzpicture}
	\draw[step=0.5cm,gray!20,very thin]  grid (17,6) ;
\end{tikzpicture}
\begin{tikzpicture}
	\draw[step=0.5cm,gray!20,very thin]  grid (17,9) ;
\end{tikzpicture}
\vskip 2.5pt 
\oppgave{}
Tegn og forklar hvordan du kobler deg til en SMART transmitter
\vskip 2.5pt 
\begin{tikzpicture}
	\draw[step=0.5cm,gray!20,very thin]  grid (17,9) ;
\end{tikzpicture}
\begin{tikzpicture}
	\draw[step=0.5cm,gray!20,very thin]  grid (17,6) ;
\end{tikzpicture}
\vskip 2.5pt 
\oppgave{}
Forklar hvordan en IO-link sensor kommuniserer med en IO-link master 
\vskip 2.5pt 
\begin{tikzpicture}
	\draw[step=0.5cm,gray!20,very thin]  grid (17,15) ;
\end{tikzpicture}
\vskip 2.5pt 
\includepdf[pages=2,angle=0]{../output/nogpl/aKom01.pdf}
\end{document}
